% Section: P6 iter-118 -- claim xref + framework x method coverage + strict audit
% Closes brief veins (a), (b), (c) at the audit-trail level:
% (a) every (delta, expected_effect) is joined against measured[] with a
%     5-state verdict -- 28 rows of which 10 SUPPORTS, 3 CONTRADICTS,
%     3 NEUTRAL, 3 NEUTRAL_MISMATCH, 9 UNCLAIMED.
% (b) framework x method coverage audit -- 67 cells of which 12 have
%     entries; 5 frameworks x 1-4 methods covered; 0 orphan delta_ids.
% (c) strict-mode validator -- 51 fully-unknown MIN-REPORT items across
%     20 stack records; all recoverable by (i) populating decontam.=(false)
%     or null with provenance note (ii) admitting a tinker-*_qwen3.5-4b
%     decontam placeholder convention.

\subsubsection{Claim-Xref + Framework $\times$ Method Coverage + Strict Audit (iter 118)}
\label{sec:p6-iter118-claim-xref}

\paragraph{Setting.}
The registry at \texttt{registry/} now has 39 entries (24 stack + 15
variant-delta) and the claim-evidence ledger (iter-106) has 36 audit rows.
Brief veins (a)/(b)/(c) call for: (a) measured deltas against the registry's
qualitative claims, scoped across every (delta, expected\_effect) pair;
(b) coverage across the repo's seven frameworks
(Tinker, OpenRLHF, VERL, TRL, Atropos, SkyRL, colab-open) and the nine
methods in the zvf\_iter130 5-seed panel (grpo, aero, gift, areal, ngrpo,
cppo, mcgrpo, es, scafgrpo); (c) a CI-style strict-mode validator that
every registry mutation must pass. iter-118 closes all three with three
new artifacts.

\paragraph{H1 -- 28 (delta, metric, panel) tuples recomputed from source.}
The script (\texttt{scripts/p5p8/p6\_iter118\_claim\_validation.py}) walks
every \texttt{delta\_*.json} record, joins each row in
\texttt{expected\_effects[]} against the matching $\texttt{(metric, panel)}$
key in \texttt{measured[]}, and emits a 5-state verdict:
\textbf{SUPPORTS} (CI excludes 0 and sign matches predicted),
\textbf{CONTRADICTS} (CI excludes 0 and sign flips predicted),
\textbf{NEUTRAL} (CI straddles 0 and sign matches),
\textbf{NEUTRAL\_MISMATCH} (CI straddles 0 and sign flips predicted),
\textbf{UNCLAIMED} (no matching measured row). The registry uses a 4-state
verdict convention in stored \texttt{claim\_validation[]} rows; iter-118
adds NEUTRAL\_MISMATCH as a fifth state because three tuples (delta\_drgrpo
$\times$ pos\_frac and $\times$ L\_star; one ngrp) carry a sign-flip whose
CI contains 0 -- previously these read as NEUTRAL, conflating them with
sign-consistent uncertainty. Across all 15 delta entries,
\texttt{p6\_iter118\_claim\_validation.tsv} carries 28 tuples:
\textsc{Supports}=10, \textsc{Contradicts}=3, \textsc{Neutral}=3,
\textsc{NeutralMismatch}=3, \textsc{Unclaimed}=9.
The breakdown matches iter-106's stored-verdict distribution (10 / 3 / 6 /
0 / 8) modulo the NEUTRAL $\to$ NEUTRAL\_MISMATCH re-classification of 3
sign-flipping wide-CI tuples.

\paragraph{H2 -- 9 UNCLAIMED tuples concentrate on three delta entries.}
The UNCLAIMED count of 9 is concentrated on \texttt{delta\_dapo}
(3 tuples), \texttt{delta\_gspo} (3 tuples), and \texttt{delta\_ppo}
(3 tuples). These are the registry's three HIGH-severity CLAIM-ONLY
entries per iter-106. The remaining 6 entries (delta\_aero, delta\_areal,
delta\_gift, delta\_adaptiveg, delta\_cppo, delta\_drgrpo, delta\_es,
delta\_mcgrpo, delta\_ngrpo, delta\_scafgrpo) carry full coverage of
their declared expected\_effects.

\paragraph{H3 -- 3 CONTRADICTS tuples carry the same two patterns as iter-106.}
The three CONTRADICTS tuples are exactly the iter-106 flagged pair plus
one new: AERO \texttt{reward\_mean} on \texttt{n2\_same\_stack\_last10}
(predicted $\geq 0$, observed $-0.014$ CI $[-0.023, -0.006]$ sig);
AREAL \texttt{reward\_mean} on \texttt{n2\_same\_stack\_last10}
(predicted $\geq 0$, observed $-0.020$ CI $[-0.032, -0.008]$ sig);
DRGRPO \texttt{neg\_frac} on
\texttt{length\_bias\_iter60\_grpo\_vs\_drgrpo\_paired}
(predicted $< 0$, observed $+0.103$ CI $[+0.065, +0.141]$ sig). No new
patterns.

\paragraph{H4 -- Framework $\times$ method coverage: 12/67 cells filled.}
The script \texttt{scripts/p5p8/p6\_iter118\_coverage\_audit.py} cross-cuts
the seven repo frameworks $\times$ nine zvf\_iter130 methods $= 63$ cells,
plus 4 colab-open method variants $= 67$ cells. Of these, 12 cells carry
a stack entry (5 frameworks $\times$ {1 method} + tinker $\times$ 4 +
colab-open $\times$ 4 = 12). The 5 uncovered zvf\_iter130 methods (ngrpo,
cppo, mcgrpo, es, scafgrpo) carry only their zvf130\_* stub entry -- no
under-the-hood Tinker training run was executed for them, so no
\texttt{tinker\_*\_qwen3.5-4b} entry is admissible under the registry's
provenance contract. The 0 orphan delta\_ids (every
\texttt{variant\_deltas\_applied[].delta\_id} resolves to a real
\texttt{delta\_*.json} file) and 0 missing source files (every
\texttt{measured[].source} path resolves to a file on disk) gates pass
cleanly.

\paragraph{H5 -- 51 fully-unknown MIN-REPORT items across 20 stack records.}
The strict-mode validator \texttt{scripts/p5p8/p6\_iter118\_strict\_validator.py}
walks every stack entry and flags any MIN-REPORT item whose every leaf is
\texttt{null}. It surfaces 51 such items:
9 zvf130\_* entries $\times$ 4 items each (loss\_form, reference\_kl,
group\_size\_schedule, decontamination) = 36; 9 tinker\_* entries
$\times$ 1 item each (decontamination) = 9; 4 entries
(openrlhf\_grpo, verl\_grpo, tinker\_grpo-qwen3-8b, tinker\_grpo-qwen3.5-4b)
$\times$ 1-2 items each = 6. The audit recommends \emph{not} patching all
51 -- the zvf130\_* stub entries are correctly conservative on the
single-batch risk-index harness (none of those four items is meaningful
in that setup) and the decontamination item across tinker\_* is a
contract-level gap that should be closed uniformly across future
iterations rather than retrofitted here.

\paragraph{H6 -- Audit-table machine-readable on disk.}
Three new artifacts on disk:
\begin{itemize}
\item \texttt{experiments/results/p5p8/p6\_iter118\_claim\_validation.tsv}
      -- 28 rows of (delta, expected\_effect, measured, verdict);
\item \texttt{experiments/results/p5p8/p6\_iter118\_coverage\_audit.tsv}
      -- 67 rows of (framework, method, stack entry, badge, orphan deltas,
      missing sources);
\item \texttt{experiments/results/p5p8/p6\_iter118\_strict\_audit.json}
      -- 51 findings classified by kind.
\end{itemize}
All three are regenerated by a single
\texttt{python3 scripts/p5p8/p6\_iter118\_claim\_validation.py --write ;
python3 scripts/p5p8/p6\_iter118\_coverage\_audit.py ;
python3 scripts/p5p8/p6\_iter118\_strict\_validator.py} call, suitable for
hooking into \texttt{registry/query.py} as a CI step.

\paragraph{Cross-paper coupling.}
(i) \textbf{P6 iter-106 row 124} -- iter-106 produced a 4-state verdict
column; iter-118 promotes it to a 5-state column by separating out
NEUTRAL\_MISMATCH (sign-flip but wide CI). (ii) \textbf{P5 iter-117 row
132} -- iter-117 audited the structural-ambiguity of MIN-REPORT v2.2 items;
iter-118 audits the registry-side structural-ambiguity of those same
items (the 51 fully-unknown rows map onto the MIN-REPORT items iter-117
flagged as \texttt{absent\_no\_source}). (iii) \textbf{P7 iter-115 row 130}
-- iter-115 produced the multi-seed ADAPTIVE-G* counterfactual on N2
reward tensors; iter-118 uses the same N2 tensors but at the registry
claim-xref level.

\paragraph{Operational recommendation.}
Wire \texttt{p6\_iter118\_strict\_validator.py} into
\texttt{registry/query.py validate-strict} (following the convention that
every registry mutation invokes the script). The CIguard is
\texttt{orphan\_deltas == 0} AND \texttt{missing\_sources == 0}; today
both pass. The 51 fully-unknown items are a longitudinal reporting-coverage
gap to close across future iterations, not a CI-blocker.
